\chapter{SYSTEM ARCHITECTURE AND METHODOLOGY}

\section{System Overview}

The Uptime Tracker system consists of three main components: a web application frontend for user interaction, a backend server for business logic and API handling, and a PostgreSQL database for persistent storage. The system also includes an automated monitoring service that periodically checks website availability.

\section{Database Architecture}

The system uses PostgreSQL database with the following schema:

\subsection{Users Table}
Stores user account information including name, email, and authentication credentials. Each user is identified by a unique ID and email address.

\subsection{Ownership Table}
Maintains relationships between websites and their owners. Includes a boolean flag for public/private visibility control. This table links website URLs to user accounts.

\subsection{Analytics Table}
Records individual ping results for monitored websites. Each entry includes:
\begin{itemize}
    \item Website URL (foreign key to ownership)
    \item Response time (ping5) in milliseconds
    \item Timestamp of the check
\end{itemize}

\subsection{Average Hour Table}
Stores aggregated hourly statistics including average response time for each website within each hour period.

\subsection{Average Day Table}
Maintains daily average statistics, providing long-term trends and historical performance data.

\section{Monitoring Methodology}

\subsection{Ping Check Process}
The monitoring service performs HTTP/HTTPS requests to registered websites every 5 minutes. For each check:
\begin{enumerate}
    \item Send HTTP request to the website URL
    \item Measure response time
    \item Record success/failure status
    \item Store result in analytics table
    \item Update hourly and daily averages
\end{enumerate}

\subsection{Data Aggregation}
Hourly and daily averages are calculated using SQL aggregate functions on the analytics data. This provides efficient querying for dashboard displays and trend analysis.

\section{User Interaction Flow}

\begin{enumerate}
    \item User registers and creates an account
    \item User adds website URLs to monitor
    \item System begins automated monitoring
    \item User views real-time status on dashboard
    \item User accesses historical analytics and reports
\end{enumerate}

\section{Technology Stack}

\begin{itemize}
    \item \textbf{Database:} PostgreSQL 18.1
    \item \textbf{Backend:} (To be determined - Python/Node.js)
    \item \textbf{Frontend:} (To be determined - React/Vue.js)
    \item \textbf{Monitoring:} Scheduled cron jobs or background workers
\end{itemize}
